\documentclass{article}
\usepackage{amsmath}  % Include the amsmath package for advanced math features
\begin{document}

The symbol you're describing, \( \binom{n}{2} \), is known as a \textbf{binomial coefficient}. It represents the number of ways to choose 2 items from a set of \( n \) items without regard to the order of selection. Mathematically, it's calculated as:

\[
\binom{n}{2} = \frac{n!}{2!(n-2)!}
\]

where \( n! \) (n factorial) is the product of all positive integers up to \( n \). For example, \( \binom{5}{2} \) would be:

\[
\binom{5}{2} = \frac{5!}{2!(5-2)!} = \frac{5 \times 4 \times 3 \times 2 \times 1}{2 \times 1 \times 3 \times 2 \times 1} = 10
\]

This means there are 10 ways to choose 2 items from a set of 5.

\end{document}
